\documentclass[11pt]{article}

\usepackage[T1]{fontenc}
\usepackage[USenglish]{babel}
\usepackage[a4paper,margin=32mm,headheight=14pt]{geometry}
\usepackage{lmodern}
\usepackage{microtype}
\usepackage{mathtools,amssymb,amsthm}
\usepackage{enumitem}
\usepackage{booktabs}
\usepackage{xcolor}
\usepackage{fancyhdr}
\usepackage[
  colorlinks=true,
  linkcolor=blue!45!black,
  citecolor=green!35!black,
  urlcolor=blue!55!black,
  pdfauthor={Alex Junior Gomez Saltachin},
  pdftitle={Arithmetic Progressions in Gap Sets of Numerical Semigroups},
  pdfsubject={Numerical semigroups and arithmetic progressions},
  pdfkeywords={numerical semigroups, gap sets, arithmetic progressions, Ap\'ery sets, Kunz coordinates, prime gaps}
]{hyperref}
\usepackage{csquotes}
\usepackage[
  backend=bibtex,
  style=numeric,
  sorting=nyt,
  giveninits=true,
  maxbibnames=99
]{biblatex}
\addbibresource{references.bib}

\theoremstyle{plain}
\newtheorem{theorem}{Theorem}[section]
\newtheorem{proposition}[theorem]{Proposition}
\newtheorem{lemma}[theorem]{Lemma}
\newtheorem{corollary}[theorem]{Corollary}
\newtheorem{conjecture}[theorem]{Conjecture}
\newtheorem{computertheorem}[theorem]{Computer-assisted theorem}
\newtheorem{computercorollary}[theorem]{Computer-assisted corollary}
\newtheorem{conditionaltheorem}[theorem]{Conditional theorem}
\theoremstyle{definition}
\newtheorem{definition}[theorem]{Definition}
\newtheorem{hypothesis}[theorem]{Hypothesis}
\newtheorem{problem}[theorem]{Problem}
\theoremstyle{remark}
\newtheorem{remark}[theorem]{Remark}

\newcommand{\N}{\mathbb N}
\newcommand{\Z}{\mathbb Z}
\newcommand{\Ap}{\operatorname{Ap}}
\newcommand{\AP}{\operatorname{AP}}
\newcommand{\compdir}{{\small\ttfamily computational\_\allowbreak
  reproducibility\_\allowbreak files}}
\newcommand{\draftstamp}{\textcolor{red!65!black}{\textsc{Draft}}}

\pagestyle{fancy}
\fancyhf{}
\fancyhead[L]{\small A. J. Gomez Saltachin}
\fancyhead[R]{\small Arithmetic progressions in gap sets}
\fancyfoot[C]{\thepage}
\fancyfoot[R]{\small\draftstamp}
\renewcommand{\headrulewidth}{0.4pt}
\renewcommand{\footrulewidth}{0pt}
\fancypagestyle{plain}{%
  \fancyhf{}
  \fancyfoot[C]{\thepage}
  \fancyfoot[R]{\small\draftstamp}
  \renewcommand{\headrulewidth}{0pt}
  \renewcommand{\footrulewidth}{0pt}}

\setlength{\parindent}{1.5em}
\setlength{\parskip}{0pt}
\setlist{topsep=0.5\baselineskip,itemsep=0.15\baselineskip,
  parsep=0pt,partopsep=0pt}
\emergencystretch=2em

\title{Arithmetic Progressions in Gap Sets of Numerical Semigroups}
\author{Alex Junior Gomez Saltachin\\[0.5em]
\small Departamento de Matem\'atica, Pontif\'{\i}cia Universidade Cat\'olica do Rio de Janeiro\\
\small Rua Marqu\^es de S\~ao Vicente, 225, Rio de Janeiro, RJ 22451-900, Brazil\\
\small \href{mailto:alex.gomez@pucp.edu.pe}{alex.gomez@pucp.edu.pe}}
\date{\small\draftstamp{} --- September 2, 2026}

\begin{document}

\maketitle

\begin{abstract}
For a numerical semigroup $S$, we study arithmetic progressions contained
in its gap set and let $\AP(S)$ denote their maximal length.  Standard
Ap\'ery theory gives $\AP_d(S)=\lfloor F/d\rfloor+1$ when the common
difference $d$ belongs to $S$, where $F$ is the Frobenius number.  This
isolates differences in the gap set, which lead to a divisor-orbit Run Lemma
in Kunz coordinates whose positional term yields an anchoring phenomenon.  For
the extremal function
\[
  M(A)=\max\{g(S):\AP(S)\leq A\},
\]
we prove $M(3)=7$ and $M(9)=54$, and computer-assisted enumeration gives
$M(8)=48$.  A prime-multiplicity construction gives the lower bound
$A(P(A+1)-1)$, where $P(x)$ is the largest prime not exceeding $x$, while
composite multiplicities incur explicit divisor-orbit losses.  This leads to
the prime-staircase conjecture
$M(A)=A(P(A+1)-1)$ for $A\neq3$.

The prime-staircase conjecture remains open in full.  We verify it
computer-assistively for $A\leq4\cdot10^{18}$, $A\neq3$, prove it
unconditionally on a set of integers $A$ of natural density one, and prove
its eventual validity under a square-root-scale hypothesis on consecutive
prime gaps.  Known prime-gap estimates also yield
$M(A)=(1+o(1))A^2$.  Finally, we obtain genus-realization results below the
prime-staircase lower bound and corresponding information about the inverse
extremal function.  Accordingly, known information on consecutive prime gaps
provides arithmetic input for the semigroup extremal problem.  The principal
remaining obstacle is structural: propagating divisor-orbit loss through the
Kunz inequalities for the residual composite multiplicities.
\end{abstract}

\medskip
\noindent\textbf{2020 Mathematics Subject Classification.}
Primary 20M14; Secondary 11B25, 11D07, 11N05.

\smallskip
\noindent\textbf{Keywords.}
Numerical semigroups, gap sets, arithmetic progressions, Ap\'ery sets,
Kunz coordinates, prime gaps.

\section{Introduction}\label{sec:introduction}

A numerical semigroup is a cofinite additive submonoid $S\subseteq\N$, where
$0\in\N$.  Its finite complement $G(S)=\N\setminus S$ is the gap set
\cite{rosales-garcia-sanchez}.
Classical invariants measure the number and location of these gaps: the
genus $g=|G(S)|$, the Frobenius number $F=\max G(S)$, the multiplicity
$m=\min(S\setminus\{0\})$, and the Ap\'ery sets.  In this paper we study a
different aspect of the gap set, namely additive patterns internal to it.
For each $d\geq1$, we define
\[
  \AP_d(S)=\max\{k:\text{some }x>0\text{ satisfies }
  x,x+d,\ldots,x+(k-1)d\in G(S)\}.
\]
We also set
\[
  \AP(S)=\max_{d\geq1}\AP_d(S),
  \qquad
  \Sigma(S)=\{(d,k)\in\N_{>0}^2:\AP_d(S)\geq k\}.
\]
The internal combinatorics of numerical-semigroup gap sets has been studied
from several complementary perspectives.  Eliahou and Fromentin's gapset
program treats $G(S)$ as a combinatorial object in its own right, using
gapsets, filtrations, multiplicity, and related invariants
\cite{eliahou-fromentin-gapsets,eliahou-fromentin-small}.  The consecutive-
difference case belongs to the classical theory of gap blocks, also called
deserts: $\AP_1(S)$ is the length of the longest such block
\cite{moree-gap-blocks}.  Residue-class statistics of gaps are naturally
expressed through Ap\'ery sets: Shor
studies equidistribution modulo an integer and, more recently, parity
distributions for free numerical semigroups
\cite{shor-equidistribution,shor-parity}.  Kunz coordinates also admit a rich
polyhedral interpretation through the Kunz cone, Kunz posets, and the Kunz
fan \cite{kaplan-oneill-group-cone,brower-mcdonough-oneill-kunz-fan}, with
related genus-counting applications \cite{bernardini-kunz}.

Rather than counting gaps or studying only their distribution among residue
classes, we study arithmetic progressions contained in the gap set itself.
Shor's questions concern residue statistics; our residue-prefix results lie
in the same Ap\'ery-set framework, but we ask for the longest ordered
patterns supported by those classes and, for divisors of the multiplicity,
by their cyclic interaction.  The special case $d=1$, corresponding to
maximal blocks of consecutive gaps, is classical; in particular, its maximum
length is $m-1$ \cite{moree-gap-blocks}.  To our knowledge, the
common-difference invariants $\AP_d(S)$ for arbitrary $d$, the global
invariant $\AP(S)$, and the extremal function $M(A)$ have not previously
been studied in this form.

The analysis rests on the following structural dichotomy:
\[
  d\in S
  \qquad\text{versus}\qquad
  d\in G(S).
\]

Subtraction by an element of $S$ preserves gaps while the result remains
positive.  Consequently, Theorem~\ref{thm:semigroup-difference} shows that,
for $d\in S_{>0}$, the Ap\'ery classes determine the entire distribution of
gaps modulo $d$ and give
\begin{equation}\label{eq:intro-semigroup-difference}
  \AP_d(S)=\left\lfloor\frac Fd\right\rfloor+1.
\end{equation}
This identity is a standard consequence of Ap\'ery-set theory.
Together with the classical consecutive-gap case, it gives the hierarchy
\[
\begin{array}{rcll}
 d=1&:&\AP_1(S)=m-1,
   &\text{classical gap blocks or deserts};\\
 d\in S_{>0}&:&\displaystyle
   \AP_d(S)=\left\lfloor\frac Fd\right\rfloor+1,
   &\text{standard Ap\'ery structure};\\
 d\in G(S)\setminus\{1\}&:&
   \multicolumn{2}{l}{\text{the genuinely difficult part of the AP spectrum}.}
\end{array}
\]
We record the semigroup-difference formulation because it isolates this last
component.

The global invariant already satisfies a useful universal bound:
\begin{equation}\label{eq:intro-universal-bound}
  \AP(S)\geq
  \max\left\{m-1,\left\lceil\frac g{m-1}\right\rceil\right\}
  \geq\lceil\sqrt g\rceil.
\end{equation}

The main semigroup mechanism is the divisor-orbit formula of
Lemma~\ref{lem:run}.  In Kunz coordinates, we write
\[
  \Ap(S,m)=\{0,w_1,\ldots,w_{m-1}\},
  \qquad w_i=i+mk_i.
\]
From the polyhedral viewpoint, the Run Lemma imposes divisor-orbit
constraints on Kunz vectors; its positional term retains information not
visible in coordinate-wise height bounds.

For $d\mid m$, a nonzero orbit under addition by $d$ reads the gap-prefixes
of its residue classes cyclically.  If its length is $\ell=m/d$, its minimum
Kunz height is $\mu$, and the first minimum occurs at zero-indexed position
$\rho$, then the longest gap run on that orbit has length
\begin{equation}\label{eq:intro-run}
  \ell\mu+\rho.
\end{equation}
The term $\rho$ is essential.  Under an upper bound on $\AP(S)$, an extremal
value of $\mu$ can force the minimum to occur at a prescribed position.  We
refer to this positional restriction as \emph{anchoring}.  The Run Lemma and
anchoring are the source of both the small exact computations and the
general penalty for composite multiplicities.

These constraints motivate the extremal function and its inverse problem,
which we define by
\[
  M(A)=\max\{g(S):\AP(S)\leq A\},
  \qquad
  a(g)=\min\{\AP(S):g(S)=g\}.
\]
The estimate \eqref{eq:intro-universal-bound} gives $M(A)\leq A^2$, but it
does not detect the arithmetic restrictions imposed by differences in
$G(S)$.

The first nontrivial extremal values are
\[
  M(3)=7,
  \qquad M(8)=48,
  \qquad M(9)=54.
\]
Here $M(8)=48$ is presently computer-assisted, while
Theorem~\ref{thm:plateau-nine} proves $M(9)=54$ by combining divisor-orbit
constraints, Kunz propagation, and a finite Kunz certificate in the remaining
multiplicities.  An explicit multiplicity-$9$
Kunz vector
\[
  (3,3,6,6,9,9,9,9)
\]
attains genus $54$ and has $\AP(S)=9$.  This is the first plateau where the
bound $M(A)\leq A^2$ is far from the true value.

The general construction is controlled by primes.  Let
\[
  P(x)=\max\{p\leq x:p\text{ is prime}\}.
\]
For a prime $p\leq A+1$, the uniform multiplicity-$p$ Kunz vector
$(A,\ldots,A)$ has genus $A(p-1)$ and AP length $A$.  This construction gives
\begin{equation}\label{eq:intro-staircase-lower}
  M(A)\geq A(P(A+1)-1).
\end{equation}
The resulting \emph{prime staircase conjecture} predicts equality for every
$A\neq 3$, while $M(3)=7$ is exceptional.  No uniqueness assertion is part of
the conjecture.  Indeed, at $A=9$ the maximum is attained by non-isomorphic
semigroups of multiplicities $7$ and $9$.  Determining the structure and
uniqueness of maximizers is therefore a separate problem from determining
the value of $M(A)$.

The opposite force comes from compositeness.  If a prime $q$ divides the
multiplicity $m$, the Run Lemma gives the genus loss
\begin{equation}\label{eq:intro-composite-loss}
  g\leq A(m-1)-\left(\frac mq-1\right)
  \left(A-\left\lfloor\frac Aq\right\rfloor\right).
\end{equation}
For the least prime factor this loss is greater than $A(\sqrt m-2)$.  Comparing
it with the extra capacity obtained by taking $m>P(A+1)$ is precisely where
estimates for gaps between consecutive primes enter the semigroup problem.

The extremal problem is therefore governed by a competition between two
effects: prime multiplicities attain the naive Kunz capacity, whereas composite
multiplicities are subject to divisor-orbit losses.  If $p=P(A+1)$, the
backward prime distance $A+1-p$ measures how much additional multiplicity a
composite competitor can have beyond the largest available prime
multiplicity.  This is where results on gaps between consecutive primes
enter the argument.

The number-theoretic input is entirely one-directional: we use known
estimates for prime gaps to obtain consequences for the semigroup extremal
problem, and we obtain no new bounds or other results concerning prime gaps.
The relevant input concerns upper bounds for gaps between consecutive
primes, or equivalently primes in short intervals, rather than the
bounded-gap problem.  Pointwise, we use the Baker--Harman--Pintz exponent
$0.525$.  For density statements we use Stadlmann's mean-square estimate for
consecutive prime gaps; Heath-Brown's classical estimate gives a weaker
historical comparison
\cite{baker-harman-pintz,heath-brown-iii,stadlmann-mean-square}.
We also record conditional consequences of standard conjectural prime-gap
bounds.

The full prime-staircase conjecture remains open.  The partial results have
three distinct scopes: the conjectural formula is verified computer-assistively
for $1\leq A\leq4\cdot10^{18}$ with $A\neq3$; it holds unconditionally for
all but $O_\varepsilon(x^{73/100+\varepsilon})$ integers $A\leq x$; and,
under the square-root gap hypothesis $g_p\leq\sqrt p-1$, it holds for every
sufficiently large $A$.

The principal unresolved step is semigroup-theoretic rather than a further
application of the prime-gap machinery.  Known prime-gap estimates allow the
criterion above to eliminate most multiplicities, but the surviving
prime-square and balanced-semiprime cases require additional loss propagated
through Run-Lemma anchoring and Kunz subadditivity.  Establishing such a
propagated-loss theorem, formulated as Problem~\ref{prob:propagated-loss}, is
the main structural route toward an unconditional proof of the full
conjecture.

The prime construction also fills the genera below its upper endpoint.  For
every $A\geq1$, Theorem~\ref{thm:realization-staircase} realizes each genus
$g$ satisfying
\[
  1\leq g\leq A(P(A+1)-1)
\]
by a numerical semigroup with $\AP(S)\leq A$.  Every proved exact jump of the
prime staircase therefore yields a corresponding plateau for $a(g)$.

Two further unconditional consequences sharpen the asymptotic picture.  The
first characterizes the square maximum:
\[
  M(A)=A^2\quad\Longleftrightarrow\quad A+1\text{ is prime}.
\]
The second concerns the deficit $\Delta(A)=A^2-M(A)$ and gives
\[
  \Delta(A)\ll A^{1.525},
  \qquad M(A)=(1+o(1))A^2.
\]
Wherever the exact staircase formula is known, the deficit satisfies
\[
  \frac{\Delta(A)}A=A+1-P(A+1),
\]
so the normalized extremal deficit is precisely the backward prime distance.

Section~\ref{sec:gap-aps} develops the elementary and Ap\'ery-theoretic
structure.  Section~\ref{sec:kunz-orbits} proves the Run Lemma and the
composite-loss theorem.  Section~\ref{sec:small-values} treats the small
extremal values and the plateau at nine.  Sections~\ref{sec:prime-staircase}
and \ref{sec:prime-gaps} develop the prime staircase and derive consequences
from known estimates for gaps between consecutive primes.
Section~\ref{sec:realization} proves the realization and
inverse-plateau results.  Section~\ref{sec:further} concludes with the
remaining structural problems, and the appendix records the computational
certificates.

\section{Arithmetic progressions in numerical semigroup gaps}
\label{sec:gap-aps}

Unless stated otherwise, $S$ has positive genus.  We use the notation
\[
  G(S)=\N\setminus S,
  \qquad g=|G(S)|,
  \qquad F=\max G(S),
  \qquad m=\min(S\setminus\{0\}).
\]
When $S=\N$, the gap set is empty; one may set $\AP_d(S)=\AP(S)=0$, but
$F$ is undefined and statements involving it do not apply.

\begin{lemma}[Downward closure]\label{lem:downward-closure}
If $x\in G(S)$ and $s\in S$ satisfy $x-s>0$, then $x-s\in G(S)$.
\end{lemma}

\begin{proof}
Otherwise $x-s\in S$, and closure under addition gives
$x=(x-s)+s\in S$, a contradiction.
\end{proof}

\begin{proposition}[Consecutive gaps]\label{prop:consecutive-gaps}
For every numerical semigroup of positive genus and multiplicity $m$, we have
\[
  \boxed{\AP_1(S)=m-1.}
\]
\end{proposition}

\begin{proof}
The integers $1,2,\ldots,m-1$ are gaps, so $\AP_1(S)\geq m-1$.  On the
other hand, any $m$ consecutive integers contain a multiple of $m$.  Since
$m\in S$ and $S$ is closed under addition, every nonnegative multiple of
$m$ belongs to $S$.  Hence no $m$ consecutive positive integers can all be
gaps, and therefore $\AP_1(S)\leq m-1$.
\end{proof}

This is the classical longest-gap-block observation; Moree uses the term
\emph{gap block} and notes the alternative terminology \emph{desert}
\cite{moree-gap-blocks}.

The consecutive case leads naturally to the more restrictive question of
when the entire gap set forms a single arithmetic progression.  There are
only two possibilities.

\begin{proposition}\label{prop:gap-set-ap}
Assume $g\geq1$.  The gap set $G(S)$ is a single arithmetic progression if
and only if either
\begin{enumerate}
  \item $S=\{0,g+1,g+2,\ldots\}$ and $G(S)=\{1,2,\ldots,g\}$; or
  \item $S=\langle2,2g+1\rangle$ and
  $G(S)=\{1,3,5,\ldots,2g-1\}$.
\end{enumerate}
For $g=1$, the two descriptions coincide.
\end{proposition}

\begin{proof}
Since $1\in G(S)$, an arithmetic progression equal to $G(S)$ has the form
\[
  \{1,1+d,\ldots,1+(g-1)d\}.
\]
If $d=1$, the semigroup is ordinary.  If $d=2$, its gaps are precisely the
odd integers below $2g+1$, and the semigroup is $\langle2,2g+1\rangle$.
If $d\geq3$, then $2,3\in S$, and therefore
$\langle2,3\rangle=\N\setminus\{1\}\subseteq S$.  It follows that
$G(S)=\{1\}$,
which is the common genus-one case.  The converse follows directly from the
two displayed gap sets.
\end{proof}

The genus-zero semigroup $\N$ has empty gap set and is excluded from this
classification.

For arbitrary common differences, the relevant structure is the distribution
of the gaps among residue classes.  Fix $d\in S_{>0}$ and recall that
\[
  \Ap(S,d)=\{s\in S:s-d\notin S\}
  =\{w_0,w_1,\ldots,w_{d-1}\},
\]
where $w_r\equiv r\pmod d$ and $0\leq r<d$.  In particular, $w_0=0$.

\begin{proposition}\label{prop:apery-residue-gaps}
For $1\leq r\leq d-1$, the gaps in residue class $r$ are
\[
  G(S)\cap(r+d\Z)=\{r,r+d,\ldots,w_r-d\}.
\]
Consequently, this residue class contains
\[
  \frac{w_r-r}{d}=\left\lfloor\frac{w_r}{d}\right\rfloor,
\]
gaps.  Taking the maximum over the nonzero residue classes gives
\[
  \AP_d(S)=\max_{w\in\Ap(S,d)}\left\lfloor\frac wd\right\rfloor.
\]
\end{proposition}

\begin{proof}
The element $w_r$ is the least member of $S$ congruent to $r$ modulo $d$.
It follows that the positive integers $r,r+d,\ldots,w_r-d$ are gaps, while
$w_r+jd\in S$ for $j\geq0$.  Writing $w_r=r+a_r d$ shows that the displayed
progression has $a_r=(w_r-r)/d=\lfloor w_r/d\rfloor$ terms.  Taking the
largest residue-class prefix gives the final formula.
\end{proof}

Summing these prefix lengths gives Selmer's genus formula
\begin{equation}\label{eq:selmer-genus}
  g=\frac1d\sum_{w\in\Ap(S,d)}w-\frac{d-1}{2}.
\end{equation}
See \cite{selmer1977}; standard background on numerical semigroups and
Ap\'ery sets is available in \cite{rosales-garcia-sanchez}.

This is the Ap\'ery residue-class framework discussed in the Introduction;
see \cite{shor-equidistribution,shor-parity} for related distributional
questions.  Here the prefix lengths will control arithmetic progressions.

When the common difference itself belongs to the semigroup, the
residue-class formula depends only on $F$ and $d$.

\begin{theorem}[Semigroup differences]\label{thm:semigroup-difference}
For every $d\in S_{>0}$, we have
\[
  \boxed{\AP_d(S)=\left\lfloor\frac Fd\right\rfloor+1.}
\]
\end{theorem}

\begin{proof}
The standard identity $\max\Ap(S,d)=F+d$ and
Proposition~\ref{prop:apery-residue-gaps} give
\[
  \AP_d(S)=\left\lfloor\frac{F+d}{d}\right\rfloor
  =\left\lfloor\frac Fd\right\rfloor+1.
\]
For a direct description, write $F=cd+r$ with $0\leq r<d$.  Since $d\in S$
and $S$ is closed under addition, every positive multiple of $d$ belongs to
$S$.  As $F\notin S$, one cannot have $d\mid F$, and hence
$1\leq r\leq d-1$.
Repeated downward closure gives the gap progression
\[
  r,r+d,\ldots,r+cd=F
\]
of length $c+1$.  Conversely, if
$x,x+d,\ldots,x+(k-1)d$ is a gap progression, then $x\geq1$ and its last
term is at most $F$.  These inequalities give
\[
  1+(k-1)d\leq x+(k-1)d\leq F,
\]
so $k\leq\lfloor F/d\rfloor+1$.
\end{proof}

\begin{corollary}[Semigroup-difference spectrum]\label{cor:s-spectrum}
The semigroup-difference part of the spectrum is
\[
  \Sigma(S)\cap(S_{>0}\times\N_{>0})
  =\{(d,k):d\in S,\ (k-1)d\leq F\}.
\]
\end{corollary}

\begin{remark}
The Ap\'ery set records the full residue-class distribution of gaps, whereas
Theorem~\ref{thm:semigroup-difference} shows that its longest prefix depends
only on $F$ and $d$.  Therefore, the part of $\Sigma(S)$ with $d\in S$ is
governed
by classical data; apart from the classical case $d=1$, the new component
occurs for $d\in G(S)\setminus\{1\}$.
\end{remark}

Although gap differences contain the new part of the spectrum, the classical
residue prefixes already impose a useful lower bound on the global invariant.

\begin{theorem}\label{thm:multiplicity-bound}
For every numerical semigroup of positive genus, we have
\[
  \boxed{\AP(S)\geq
  \max\left\{m-1,\left\lceil\frac{g}{m-1}\right\rceil\right\}.}
\]
\end{theorem}

\begin{proof}
Proposition~\ref{prop:consecutive-gaps} gives $\AP_1(S)=m-1$.  By
Proposition~\ref{prop:apery-residue-gaps} at $d=m$, the gaps split into
$m-1$ nonzero residue classes modulo $m$.  Writing
$k_i=(w_i-i)/m$ for their lengths gives $\sum_i k_i=g$ and
\[
  \AP_m(S)=\max_i k_i
  \geq\left\lceil\frac g{m-1}\right\rceil.
\]
Combining these two bounds gives
\[
  \AP(S)\geq\max\{\AP_1(S),\AP_m(S)\}
  \geq\max\left\{m-1,
  \left\lceil\frac g{m-1}\right\rceil\right\}.
\]
\end{proof}

\begin{corollary}\label{cor:sqrt-bound}
Every numerical semigroup of positive genus satisfies
\[
  \boxed{\AP(S)\geq\lceil\sqrt g\rceil.}
\]
\end{corollary}

\begin{proof}
Apply $\max\{x,g/x\}\geq\sqrt g$ with $x=m-1$ and use integrality.
\end{proof}

Combining Theorems~\ref{thm:semigroup-difference} and
\ref{thm:multiplicity-bound} at $d=m$ gives
\[
  \AP_m(S)=\left\lfloor\frac Fm\right\rfloor+1
  \geq\left\lceil\frac g{m-1}\right\rceil.
\]
Consequently, the genus satisfies
\[
  g\leq(m-1)\left(\left\lfloor\frac Fm\right\rfloor+1\right).
\]
\section{Kunz coordinates and divisor-orbit constraints}
\label{sec:kunz-orbits}

With respect to the multiplicity, we write
\[
  \Ap(S,m)=\{0,w_1,\ldots,w_{m-1}\},
  \qquad w_i=i+mk_i.
\]
The positive integers $k_i$ are the Kunz coordinates.  The coordinate $k_i$
is the number of gaps congruent to $i$ modulo $m$, so
\begin{equation}\label{eq:kunz-genus}
  g=\sum_{i=1}^{m-1}k_i.
\end{equation}
For $1\leq i,j\leq m-1$, the Kunz coordinates satisfy
\begin{align}
  k_i+k_j&\geq k_{i+j} &&(i+j<m),\label{eq:kunz-unwrapped}\\
  k_i+k_j+1&\geq k_{i+j-m} &&(i+j>m).\label{eq:kunz-wrapped}
\end{align}
When $i+j=m$, no nontrivial coordinate inequality is required.  Conversely,
every positive integral vector satisfying these inequalities is the Kunz
vector of a numerical semigroup of multiplicity $m$; see
\cite{rosales-garcia-sanchez}.  The polyhedral background appears in
\cite{kaplan-oneill-group-cone,brower-mcdonough-oneill-kunz-fan}, with
related genus-counting applications in \cite{bernardini-kunz}.

If $\AP(S)\leq A$, then Proposition~\ref{prop:apery-residue-gaps} at $d=m$
gives
\begin{equation}\label{eq:kunz-box-general}
  1\leq k_i\leq A
  \qquad(1\leq i<m).
\end{equation}

The coordinatewise box \eqref{eq:kunz-box-general} does not record how a
progression moves cyclically through residue classes.  When its common
difference divides the multiplicity, this interaction is exact.

\begin{lemma}[Run Lemma]\label{lem:run}
Let $d\mid m$ with $1\leq d<m$, and let $1\leq r\leq d-1$.  In cyclic
order, the orbit of $r$ under addition by $d$ modulo $m$ is
\[
  c_0=r,c_1=r+d,\ldots,c_{\ell-1}=r+m-d,
  \qquad \ell=\frac md.
\]
Define the minimum height and its first position by
\[
  \mu=\min_{0\leq j<\ell}k_{c_j},
  \qquad \rho=\min\{j:k_{c_j}=\mu\}.
\]
The longest gap progression of difference $d$ contained in this orbit then
has length
\[
  \boxed{\AP_d(S)|_{\mathcal O_r}=\ell\mu+\rho
  =\frac md\,\mu+\rho.}
\]
\end{lemma}

\begin{proof}
We read the positive integers in the orbit with step $d$:
\[
  r,r+d,r+2d,\ldots.
\]
If $n=h\ell+j$ with $0\leq j<\ell$, then
$r+nd=c_j+hm$.  The class $c_j$ contributes the gap-prefix
\[
  c_j,c_j+m,\ldots,c_j+(k_{c_j}-1)m.
\]
In this way, the orbit is read cyclically, one complete cycle at each height
level $h$.  At every level $h<\mu$, all $\ell$ entries are gaps, producing
$\ell\mu$ initial gaps.  At level $h=\mu$, the first $\rho$ classes still
contribute gaps, while $c_\rho+\mu m$ is the first nongap.  The initial run
has length $\ell\mu+\rho$.  At every later level, the class $c_\rho$ again
contributes a nongap, so every later run has length at most $\ell-1$.  Since
$\mu\geq1$, the initial run is strictly longer.
\end{proof}

The positional term $\rho$ strengthens the minimum-height constraint whenever
the latter is close to its largest possible value.

\begin{corollary}[Orbit constraint]\label{cor:orbit-constraint}
If $\AP(S)\leq A$, then every nonzero divisor orbit in
Lemma~\ref{lem:run} satisfies
\[
  \boxed{\ell\mu+\rho\leq A.}
\]
In particular, $\mu\leq\lfloor A/\ell\rfloor$; if equality holds, the first
minimum is restricted by
\[
  \rho\leq A-\ell\left\lfloor\frac A\ell\right\rfloor.
\]
\end{corollary}

\begin{remark}[Anchoring]
The position $\rho$ contains information absent from the minimum-height bound.
When $\mu$ is extremal, the first minimum can be forced to the beginning of
the cyclic order.  This anchoring phenomenon is decisive in the plateau
calculation of Section~\ref{sec:small-values}.
\end{remark}

Summing the orbit constraints converts this local information into a global
genus loss for composite multiplicities.

\begin{theorem}[Composite-multiplicity loss bound]
\label{thm:composite-loss}
Let $\AP(S)\leq A$, let $m$ be the multiplicity, and let $q\mid m$ be prime.
Under these assumptions, the genus satisfies
\[
  \boxed{g\leq
  \left(\frac mq-1\right)
  \left((q-1)A+\left\lfloor\frac Aq\right\rfloor\right)
  +(q-1)A.}
\]
Equivalently, the bound can be written as
\begin{equation}\label{eq:composite-loss}
  \boxed{g\leq A(m-1)-\left(\frac mq-1\right)
  \left(A-\left\lfloor\frac Aq\right\rfloor\right).}
\end{equation}
\end{theorem}

\begin{proof}
We set $d=m/q$.  Addition by $d$ modulo $m$ has $m/q$ cycles of length $q$.
One contains residue $0$; the other $m/q-1$ cycles contain only nonzero
residues.  On each nonzero cycle, Corollary~\ref{cor:orbit-constraint} gives
$\mu\leq\lfloor A/q\rfloor$.  Its coordinate sum is therefore at most
\[
  (q-1)A+\left\lfloor\frac Aq\right\rfloor.
\]
The zero cycle contains $q-1$ nonzero coordinates and contributes at most
$(q-1)A$.  Adding the contributions of all cycles gives
\[
\begin{aligned}
g
&\leq
\left(\frac mq-1\right)
\left((q-1)A+\left\lfloor\frac Aq\right\rfloor\right)
+(q-1)A\\
&=A(m-1)-\left(\frac mq-1\right)
\left(A-\left\lfloor\frac Aq\right\rfloor\right),
\end{aligned}
\]
which is \eqref{eq:composite-loss}.
\end{proof}

\begin{remark}
The term $A(m-1)$ is the naive genus capacity from
\eqref{eq:kunz-box-general}.  The subtracted term in
\eqref{eq:composite-loss} is the genus penalty forced by the divisor orbits.
\end{remark}

\section{Small extremal values and the plateau at nine}
\label{sec:small-values}

The composite-multiplicity loss already determines the exceptional value at
$A=3$.

Applying Theorem~\ref{thm:composite-loss} with $(A,m,q)=(3,4,2)$ gives
\[
  g\leq3(4-1)-(2-1)\left(3-\left\lfloor\frac32\right\rfloor\right)=7.
\]
For $m\leq3$, we have $g\leq3(m-1)\leq6$.  To attain the remaining upper
bound, consider the semigroup
\[
  S=\langle4,7,13\rangle,
\]
which has Kunz vector $(3,3,1)$ and gap set
\[
  G(S)=\{1,2,3,5,6,9,10\}.
\]
For differences $1,2,3,4$, the following progressions have length three,
respectively:
\[
  (1,2,3),\quad(1,3,5),\quad(3,6,9),\quad(1,5,9).
\]
A direct check of the displayed gap set shows that no progression with any
of these four differences has four terms.  If $d\geq5$, a four-term
progression has span $3d>F-1=9$.  We conclude that $\AP(S)=3$, and hence
\begin{equation}\label{eq:M3}
  \boxed{M(3)=7.}
\end{equation}

At $A=8$, a direct exhaustive search gives the exact extremal genus.

\begin{computertheorem}[The value at eight]\label{thm:M8-computation}
The extremal value at $A=8$ is
\[
  \boxed{M(8)=48.}
\]
\end{computertheorem}

\begin{proof}[Computer-assisted proof]
The accompanying directory \compdir{}
contains the program \path{direct_ap_search.c}, which first verifies the
prime-multiplicity
lower-bound example directly.  The naive coordinate bound then excludes
multiplicities at or below that prime, and the program performs exhaustive
branch-and-bound over the remaining multiplicities $m\leq9$ and coordinate
vectors with $1\leq k_i\leq8$.  It tests the Kunz inequalities and prunes
only when the largest possible genus on a branch cannot exceed the verified
lower bound.  For every remaining candidate, it constructs the finite gap set,
computes gap arithmetic progressions directly, and maximizes $\sum_i k_i$.
It uses no Run-Lemma constraint and returns $M(8)=48$.  The corresponding
output is designated \path{direct_ap_A8.log} in the same directory.
The maximum is attained at prime multiplicity $7$ by the construction of
Proposition~\ref{prop:prime-staircase} below.
\end{proof}

For comparison, Theorem~\ref{thm:composite-loss} with $(A,m,q)=(8,8,2)$
gives the stronger multiplicity-specific bound $g\leq44$.

At $A=9$, the single-divisor estimate alone does not recover the sharp bound.
Assume that $\AP(S)\leq9$; the Kunz coordinates then satisfy
$1\leq k_i\leq9$.

For $m=10$ and $d=5$, the relevant pairs are
$(k_r,k_{r+5})$, where $1\leq r\leq4$.  The Run Lemma gives
\begin{equation}\label{eq:m10-pairs}
  \min(k_r,k_{r+5})\leq4.
\end{equation}
For $d=2$, the odd orbit $(1,3,5,7,9)$ has length $5$.  The corresponding
orbit constraint is therefore
\begin{equation}\label{eq:m10-odd}
  \min(k_1,k_3,k_5,k_7,k_9)=1.
\end{equation}

For $m=9$ and $d=3$, the two nonzero cycles are $(1,4,7)$ and $(2,5,8)$.
On either cycle, we have $3\mu+\rho\leq9$.  This inequality implies that
$\mu\leq3$ and, if $\mu=3$, that $\rho=0$.  More explicitly, it gives
\begin{equation}\label{eq:m9-anchor}
  \begin{split}
  \min(k_1,k_4,k_7)=3&\Longrightarrow k_1=3,\\
  \min(k_2,k_5,k_8)=3&\Longrightarrow k_2=3.
  \end{split}
\end{equation}

For $m=8$ and $d=4$, the pair constraints are
\begin{equation}\label{eq:m8-pairs}
  \min(k_i,k_{i+4})\leq4\qquad(1\leq i\leq3).
\end{equation}
The $d=2$ odd orbit $(1,3,5,7)$ also satisfies
\begin{equation}\label{eq:m8-odd}
  4\mu+\rho\leq9.
\end{equation}

\begin{lemma}\label{lem:small-multiplicity}
If $\AP(S)\leq9$ and $m\leq7$, then the genus satisfies $g\leq54$.
\end{lemma}

\begin{proof}
The coordinate bound gives
$g=\sum_{i=1}^{m-1}k_i\leq9(m-1)\leq54$.
\end{proof}

\begin{proposition}\label{prop:m8-hand}
If $\AP(S)\leq9$ and $m=8$, then the genus satisfies $g\leq48$.
\end{proposition}

\begin{proof}
Theorem~\ref{thm:composite-loss} with $q=2$ gives
\[
  g\leq9(8-1)-(4-1)\left(9-\left\lfloor\frac92\right\rfloor\right)=48.
\]
Equivalently, each pair in \eqref{eq:m8-pairs} contributes at most $13$,
and $k_4\leq9$.
\end{proof}

For the remaining multiplicities, anchoring interacts with Kunz
subadditivity.  For instance, if $m=9$ and
$k_1=k_2=3$, then
\[
  k_3\leq k_1+k_2=6,
  \qquad k_4\leq2k_2=6,
\]
so $g\leq3+3+6+6+9+9+9+9=54$.  The doubling map on nonzero residues
contains the cycle
\[
  1\to2\to4\to8\to7\to5\to1\pmod9,
\]
which propagates such bounds through the remaining coordinates.

\begin{proposition}[Finite Kunz certification]\label{prop:finite-certificate}
For $m=8,9,10$, consider the maximum of $\sum_i k_i$ over vectors in
$\{1,\ldots,9\}^{m-1}$ satisfying all Kunz inequalities and the orbit
constraints \eqref{eq:m10-pairs}--\eqref{eq:m8-odd}.  Exhaustive enumeration
gives
\[
\begin{array}{ccc}
  \toprule
  m&\max\sum_i k_i&\text{unique maximizing vector}\\
  \midrule
  8&44&(7,9,2,9,4,4,9)\\
  9&54&(3,3,6,6,9,9,9,9)\\
  10&53&(9,9,9,8,6,4,4,3,1).\\
  \bottomrule
\end{array}
\]
\end{proposition}

\begin{proof}[Computer-assisted proof]
The finite search checks only $1\leq k_i\leq9$, the Kunz inequalities, and
the displayed Run-Lemma consequences.  It neither constructs gap sets nor
computes any $\AP_d(S)$.  It was performed by
the program \path{plateau_certificate.c} in the accompanying directory
\compdir{}; the corresponding output is
designated \path{plateau_certificate.log}.  An independent reordered
branch-and-bound implementation is provided in
\path{plateau_certificate_independent.c}, with output in
\path{plateau_certificate_independent.log}.  The constraints are
necessary for $\AP(S)\leq9$; no converse characterization is asserted.  They
suffice for the stated optimization.
\end{proof}

The preceding bounds cover every possible multiplicity and therefore give
the desired plateau estimate.

\begin{theorem}[Plateau at nine]\label{thm:plateau-nine}
If $\AP(S)\leq9$, then the genus satisfies
\[
  \boxed{g(S)\leq54.}
\]
\end{theorem}

\begin{proof}
Theorem~\ref{thm:multiplicity-bound} gives $m-1\leq9$, so $m\leq10$.
Lemma~\ref{lem:small-multiplicity} handles $m\leq7$,
Proposition~\ref{prop:m8-hand} handles $m=8$, and
Proposition~\ref{prop:finite-certificate} gives $g\leq54$ for $m=9$ and
$g\leq53$ for $m=10$.
\end{proof}

Consider the multiplicity-$9$ Kunz vector
\begin{equation}\label{eq:extremal-vector}
  (k_1,\ldots,k_8)=(3,3,6,6,9,9,9,9).
\end{equation}
It satisfies the Kunz inequalities.  Indeed, the minimum left sides of the
unwrapped inequalities with targets $2,\ldots,8$ are
$6,6,6,9,9,12,12$, while those of the wrapped inequalities with targets
$1,\ldots,7$ are $13,16,16,19,19,19,19$.  Its genus is $54$, and
\[
  F=\max_{1\leq i\leq8}(i+9k_i)-9=80.
\]

\begin{proposition}\label{prop:sharpness-nine}
The semigroup defined by \eqref{eq:extremal-vector} satisfies $\AP(S)=9$.
\end{proposition}

\begin{proof}
Suppose first that $\gcd(d,9)=1$.  Every nine consecutive terms of a
$d$-progression
represent all residues modulo $9$, so one is a positive multiple of $9$ and
belongs to $S$.  Therefore, $\AP_d(S)\leq8$.

For $d=3$, the nonzero height triples are $(3,6,9)$ and $(3,9,9)$.  The Run
Lemma gives length $3\cdot3=9$ on each; the zero orbit has a multiple of $9$
every third term.  It follows that $\AP_3(S)=9$.

For $d=6$, direct inspection of the residue prefixes gives the following
longest-run lengths:
\[
\begin{array}{c*{6}{c}}
  \toprule
  x\bmod6&1&2&3&4&5&0\\
  \midrule
  \text{length}&6&6&2&4&4&2.\\
  \bottomrule
\end{array}
\]
For example, the two length-six runs are
$(1,7,13,19,25,31)$ and $(2,8,14,20,26,32)$; their next terms are nongaps.
This is a direct inspection, since $6\nmid9$.

If $\gcd(d,9)=3$ and $d\notin\{3,6\}$, then $d\geq12$, and every gap lies
in $[1,80]$.  Therefore, the progression has length at most
$\lfloor79/12\rfloor+1=7$.  Finally, if $9\mid d$, the progression remains
in one residue class, containing at most $9$ gaps.  Equality occurs for
$d=9$.  These cases cover all $d>0$.
\end{proof}

Combining the upper bound with the sharpness construction gives
\begin{equation}\label{eq:M9}
  \boxed{M(9)=54.}
\end{equation}
The estimate $M(9)\leq81$ from Corollary~\ref{cor:sqrt-bound} misses the
loss imposed by gap differences; the plateau calculation exhibits the role
of anchoring and Kunz propagation.

\section{The prime staircase and composite multiplicities}
\label{sec:prime-staircase}

For $x\geq2$, we denote by $P(x)$ the largest prime not exceeding $x$.

\begin{proposition}[Prime staircase construction]
\label{prop:prime-staircase}
Let $p\leq A+1$ be prime.  The semigroup with multiplicity $p$ and uniform
Kunz vector $(A,\ldots,A)$ satisfies
\[
  \AP(S)=A,
  \qquad g=A(p-1).
\]
\end{proposition}

\begin{proof}
The uniform vector satisfies all Kunz inequalities.  If $p\nmid d$, every
$p$ consecutive terms of a $d$-progression contain a positive multiple of
$p$, so $\AP_d(S)\leq p-1\leq A$.  If $p\mid d$, the progression stays in
one residue class, and each nonzero class contains exactly $A$ gaps.  Hence
$\AP_d(S)\leq A$, with equality for $d=p$.
\end{proof}

\begin{corollary}[Prime staircase lower bound]
\label{cor:prime-staircase-lower}
For every $A\geq1$, we have
\[
  \boxed{M(A)\geq A(P(A+1)-1).}
\]
\end{corollary}

\begin{theorem}[Quadratic growth]\label{thm:quadratic-growth}
For every $A\geq1$, the extremal function satisfies
\[
  \frac{A(A-1)}2\leq M(A)\leq A^2.
\]
In particular, $M(A)=\Theta(A^2)$.
\end{theorem}

\begin{proof}
The upper bound is Corollary~\ref{cor:sqrt-bound}.  Bertrand's postulate
gives $P(A+1)>(A+1)/2$ for $A\geq2$, and the lower estimate follows from
Corollary~\ref{cor:prime-staircase-lower}; the case $A=1$ follows directly.
\end{proof}

The prime construction and the quadratic upper bound suggest that the
preceding lower bound is usually sharp.

\begin{conjecture}[Prime staircase conjecture]
\label{conj:prime-staircase}
For every $A\geq1$ with $A\neq 3$, the conjectural formula is
\[
  \boxed{M(A)=A(P(A+1)-1).}
\]
At the exceptional value, $M(3)=7$.
\end{conjecture}

Direct computation for the first values gives the following comparison.
\begingroup
\small
\[
\begin{array}{c*{12}{r}}
 \toprule
 A&1&2&3&4&5&6&7&8&9&10&11&12\\
 \midrule
 M(A)&1&4&7&16&20&36&42&48&54&100&110&144\\
 A(P(A+1)-1)&1&4&6&16&20&36&42&48&54&100&110&144.\\
 \bottomrule
\end{array}
\]
\endgroup
The rows agree for $1\leq A\leq12$ except at $A=3$, the exceptional value
established in \eqref{eq:M3}.  Therefore, from $A=4$ through $12$, the prime
staircase exactly predicts the directly computed extremal genus.  This is
computational evidence, not a proof of the conjecture, and is independent of
the divisor-orbit criterion used below.

The conjecture does not assert uniqueness.  At $A=9$, both the uniform
multiplicity-$7$ semigroup and the multiplicity-$9$ semigroup in
\eqref{eq:extremal-vector} are maximizers.

\begin{proposition}[Composite-multiplicity reduction]
\label{prop:composite-reduction}
Let $p=P(A+1)$.  To prove the conjectured upper bound, it suffices to consider
\[
  \boxed{p<m\leq A+1.}
\]
Every such multiplicity is composite.
\end{proposition}

\begin{proof}
Theorem~\ref{thm:multiplicity-bound} gives $m\leq A+1$.  If $m\leq p$, then
$g\leq A(m-1)\leq A(p-1)$.  By definition of $p$, there is no prime $r$
satisfying $p<r\leq A+1$.
\end{proof}

For a prime divisor $q\mid m$, Theorem~\ref{thm:composite-loss} proves the
desired upper bound whenever
\begin{equation}\label{eq:loss-versus-capacity}
  \left(\frac mq-1\right)
  \left(A-\left\lfloor\frac Aq\right\rfloor\right)
  \geq A(m-p).
\end{equation}
The left side is the Run-Lemma loss, and the right side is the excess capacity
obtained by using $m>p$.

To compare these two quantities uniformly, we use the least prime factor of
the competing multiplicity.

\begin{lemma}[Least-factor loss bound]\label{lem:least-factor-loss}
Let $m$ be composite and $q_0$ its least prime factor.  The associated loss
satisfies
\[
  \boxed{
  \left(\frac{m}{q_0}-1\right)
  \left(A-\left\lfloor\frac{A}{q_0}\right\rfloor\right)
  \geq A\frac{(\sqrt m-1)^2}{\sqrt m}
  >A(\sqrt m-2).}
\]
\end{lemma}

\begin{proof}
We first observe that
\[
  A-\left\lfloor\frac A{q_0}\right\rfloor
  \geq A\left(1-\frac1{q_0}\right).
\]
To analyze the remaining factor, define
\[
  h(t)=\left(\frac mt-1\right)\left(1-\frac1t\right).
\]
For this function, differentiation gives
\[
  h'(t)=\frac{2m-(m+1)t}{t^3}<0
  \qquad(t\geq2).
\]
Since $q_0\leq\sqrt m$ and $h$ is decreasing on the relevant interval, we
obtain
\[
  h(q_0)\geq h(\sqrt m)
  =\frac{(\sqrt m-1)^2}{\sqrt m}
  =\sqrt m-2+\frac1{\sqrt m}>\sqrt m-2.
\]
Multiplying by $A$ proves the result.
\end{proof}

The least-factor estimate rules out every composite competitor once the
backward distance to the preceding prime is sufficiently small.

\begin{theorem}[Prime-gap criterion]\label{thm:prime-gap-criterion}
Let $p=P(A+1)$, and assume that
\begin{equation}\label{eq:sqrt-gap-criterion}
  \boxed{A+1-p\leq\sqrt p-2}.
\end{equation}
Under this assumption, we have
\[
  \boxed{M(A)=A(p-1).}
\]
\end{theorem}

\begin{proof}
By Proposition~\ref{prop:composite-reduction}, only composite
$p<m\leq A+1$ require consideration.  Theorem~\ref{thm:composite-loss} and
Lemma~\ref{lem:least-factor-loss} give
\[
  g<A(m-1)-A(\sqrt m-2).
\]
Since $m>p$ and \eqref{eq:sqrt-gap-criterion} holds, we have
\[
  m-p\leq A+1-p\leq\sqrt p-2<\sqrt m-2.
\]
Combining these inequalities yields
$g<A(m-1)-A(m-p)=A(p-1)$.  The prime staircase construction
attains $A(p-1)$.
\end{proof}

The theorem identifies the threshold at which prime-gap estimates imply
optimality of the prime construction: the backward distance from $A+1$ to
its preceding prime must be at most the square-root threshold.  A single
divisor does not always suffice.  At $A=9$, $p=7$, and
\[
  g\leq60\quad(m=9,q=3),
  \qquad g\leq61\quad(m=10,q=2),
\]
whereas the sharp value is $54$.  The additional anchoring and Kunz
propagation of Section~\ref{sec:small-values} are genuinely stronger.

\section{Consequences of prime-gap estimates}\label{sec:prime-gaps}

Known information on consecutive prime gaps enters the extremal problem
through Theorem~\ref{thm:prime-gap-criterion}.  It yields three statements
with distinct statuses: finite computer-assisted verification, unconditional
density-one validity, and conditional eventual validity.

\begin{proposition}[Finite verification of the composite criterion]
\label{prop:criterion-computation}
For $1\leq A\leq10^8$, every composite integer $m$ satisfying
$P(A+1)<m\leq A+1$ also satisfies inequality~\eqref{eq:loss-versus-capacity}
for at least one prime divisor $q\mid m$, except precisely for
\[
  (A,m)=(3,4),(8,9),(9,9),(9,10).
\]
\end{proposition}

\begin{proof}[Computer-assisted proof]
The accompanying directory \compdir{}
contains \path{staircase_criterion.c}.  This exact integer computation factors
every relevant $m$ and tests \eqref{eq:loss-versus-capacity}.  It does not enumerate
semigroups or compute gap progressions.  The completed run checked
$1{,}449{,}077{,}567$ pairs and returned exactly the four displayed failures.
Its output is recorded in \path{staircase_criterion_1e8.log} in the same
directory.
\end{proof}

The case $(3,4)$ gives the genuine exception \eqref{eq:M3}; the case
$(8,9)$ is covered by Theorem~\ref{thm:M8-computation}; and the two cases at
$A=9$ are covered by Theorem~\ref{thm:plateau-nine}.

\begin{computertheorem}[Prime staircase formula in the verified range]
\label{thm:verified-range}
For every integer $A$ satisfying
\[
  1\leq A\leq4\cdot10^{18},
  \qquad A\neq 3,
\]
we have
\[
  \boxed{M(A)=A(P(A+1)-1).}
\]
Moreover, $M(3)=7$.
\end{computertheorem}

\begin{proof}
Proposition~\ref{prop:criterion-computation} and the three special cases
settle $A\leq10^8$.  Above that range, the computation of Oliveira e Silva,
Herzog, and Pardi shows that prime gaps with lower prime below
$4\cdot10^{18}$ are at most $1476$
\cite{oliveira-silva-herzog-pardi,oliveira-gap-table}.  Therefore, for composite
$A+1$ in the remaining interval, the backward prime distance satisfies
\[
  A+1-P(A+1)\leq1475<\sqrt{P(A+1)}-2,
\]
and Theorem~\ref{thm:prime-gap-criterion} applies.  If $A+1$ is prime, the
left side of \eqref{eq:sqrt-gap-criterion} is zero, so the hypothesis also
holds.
\end{proof}

The finite range does not control the possible exceptions beyond its
endpoint.  To obtain an unconditional asymptotic statement, let $p^+$ denote
the prime following $p$ and define $g_p=p^+-p$.

\begin{theorem}[Density-one prime staircase]\label{thm:density-one}
For every $\varepsilon>0$, the exceptional set satisfies
\[
  \#\{A\leq x:M(A)\neq A(P(A+1)-1)\}
  =O_\varepsilon(x^{73/100+\varepsilon}).
\]
In particular, the prime staircase formula holds on a set of natural density
$1$.
\end{theorem}

\begin{proof}
If \eqref{eq:sqrt-gap-criterion} fails for $p=P(A+1)$, then, apart from
finitely many small primes, $g_p>\sqrt p-1$.  At most $g_p$ potentially
exceptional values of $A$ lie in the gap following $p$.  This observation
gives
\[
  E(x)\ll\sum_{\substack{p\leq x+1\\g_p>\sqrt p-1}}g_p.
\]
On a dyadic block $X<p\leq2X$, and for $X$ sufficiently large,
$\sqrt p-1\geq\sqrt X-1\geq\frac12\sqrt X$.  Consequently, on the
restricted set where $g_p>\sqrt p-1$, we have
\[
  g_p\leq2X^{-1/2}g_p^2.
\]
It follows that
\[
  \sum_{\substack{X<p\leq2X\\g_p>\sqrt p-1}}g_p
  \ll X^{-1/2}\sum_{X<p\leq2X}g_p^2.
\]
Stadlmann's estimate
\[
  \sum_{p\leq Y}g_p^2
  \ll_\varepsilon Y^{123/100+\varepsilon}
\]
therefore gives $O_\varepsilon(X^{73/100+\varepsilon})$ on each block
\cite{stadlmann-mean-square}.  Dyadic summation proves the theorem.
\end{proof}

This is a density-one conclusion: it allows a zero-density exceptional set
and does not assert that the formula holds for every sufficiently large $A$.

\begin{remark}
If one instead uses Heath-Brown's classical mean-square estimate
\[
  \sum_{p\leq X}g_p^2
  \ll_\varepsilon X^{23/18+\varepsilon},
\]
then the same argument gives the weaker exceptional-set bound
$O_\varepsilon(x^{7/9+\varepsilon})$ \cite{heath-brown-iii}.  Stadlmann's
more recent improvement supplies the sharper density-one estimate stated
above.
\end{remark}

The density-one theorem still permits infinitely many exceptions.  Eventual
validity follows from the following square-root-scale hypothesis.

\begin{hypothesis}[Square-root gap hypothesis $H_{\sqrt{\ }}$]
\label{hyp:sqrt-gap}
For all sufficiently large primes $p$, we have
\[
  g_p\leq\sqrt p-1.
\]
\end{hypothesis}

\begin{conditionaltheorem}\label{thm:conditional-staircase}
Under Hypothesis~\ref{hyp:sqrt-gap}, the formula
\[
  \boxed{M(A)=A(P(A+1)-1)}
\]
holds for all sufficiently large $A$.
\end{conditionaltheorem}

\begin{proof}
We set $p=P(A+1)$.  If $A+1$ is composite, then
$A+1-p\leq g_p-1\leq\sqrt p-2$; if it is prime, the left side is zero.
Theorem~\ref{thm:prime-gap-criterion} applies.
\end{proof}

Several standard conjectures on prime gaps are much stronger than
Hypothesis~\ref{hyp:sqrt-gap}.  For example, Cram\'er's conjectural bound
$g_p=O((\log p)^2)$, as well as the usual prime-gap consequences of
Firoozbakht's conjecture, imply Hypothesis~\ref{hyp:sqrt-gap} for all
sufficiently large $p$
\cite{cramer-prime-gaps,kourbatov-firoozbakht}.  Therefore, either conjecture
would imply the eventual prime-staircase formula through
Theorem~\ref{thm:prime-gap-criterion}.  Andrica's conjecture, by contrast,
does not provide the constant required by Hypothesis~\ref{hyp:sqrt-gap}
\cite{andrica-conjecture}.

Independently of any prime-gap hypothesis, the quadratic upper bound is
attained exactly when $A+1$ is prime.

\begin{theorem}\label{thm:square-characterization}
For every $A\geq1$, we have
\[
  \boxed{M(A)=A^2\quad\Longleftrightarrow\quad A+1\text{ is prime}.}
\]
\end{theorem}

\begin{proof}
If $A+1$ is prime, Proposition~\ref{prop:prime-staircase} gives genus $A^2$,
and Corollary~\ref{cor:sqrt-bound} gives the reverse inequality.  Suppose
$A+1$ is composite.  If $m\leq A$, then
$g\leq A(m-1)\leq A(A-1)<A^2$.  If $m=A+1$, any prime divisor $q\mid m$
gives a strictly positive loss in \eqref{eq:composite-loss}.  These exhaust
the possible multiplicities.
\end{proof}

Even when the exact staircase formula is not known, the distance to the
preceding prime controls its deficit.  We define
\[
  \Delta(A)=A^2-M(A).
\]

\begin{theorem}[Unconditional deficit estimate]\label{thm:deficit}
For every $A\geq1$, the deficit satisfies
\[
  \boxed{0\leq\Delta(A)\leq A(A+1-P(A+1)).}
\]
In addition, the pointwise estimate gives
\[
  \boxed{\Delta(A)\ll A^{1.525}},
  \qquad
  \boxed{M(A)=A^2+O(A^{1.525})=(1+o(1))A^2}.
\]
\end{theorem}

\begin{proof}
The first assertion combines $M(A)\leq A^2$ with
Corollary~\ref{cor:prime-staircase-lower}.  Baker, Harman, and Pintz prove
\[
  A+1-P(A+1)\ll A^{0.525}
\]
for sufficiently large $A$ \cite{baker-harman-pintz}.  This is the relevant
unconditional pointwise short-interval estimate.  Substitution gives the
deficit estimate, and division by $A^2$ gives the asymptotic equivalence.
\end{proof}

At every $A\neq 3$ for which the prime staircase formula is proved, the
normalized deficit is
\begin{equation}\label{eq:normalized-deficit}
  \boxed{\frac{\Delta(A)}A=A+1-P(A+1).}
\end{equation}
Therefore, throughout the computer-verified range and on the density-one set of
Theorem~\ref{thm:density-one}, the normalized deficit is exactly the backward
prime distance.  Outside the proved cases, \eqref{eq:normalized-deficit}
remains conjectural.

\section{Genus realization and inverse plateaux}\label{sec:realization}

The prime construction can be varied without losing control of the AP
length.  Allowing the Kunz coordinates to range in a bounded band produces
all intermediate genera needed below.

\begin{lemma}[Bounded-band realization]\label{lem:bounded-band}
Let $p\leq A+1$ be prime, let $1\leq B\leq A$, and set
$L_B=\lceil B/2\rceil$.  Every vector
\[
  (k_1,\ldots,k_{p-1})\in\{L_B,L_B+1,\ldots,B\}^{p-1}
\]
is a valid Kunz vector of multiplicity $p$, and its semigroup satisfies
\[
  \boxed{\AP(S)\leq\max\{p-1,B\}\leq A.}
\]
\end{lemma}

\begin{proof}
For every target coordinate, we have
$k_i+k_j\geq2L_B\geq B\geq k_r$.  Hence all unwrapped Kunz inequalities hold;
the wrapped inequalities are weaker.  If $p\nmid d$, every $p$ consecutive
terms contain a positive multiple of $p$, giving $\AP_d(S)\leq p-1$.  If
$p\mid d$, a progression remains in one residue class, which contains at
most $B$ gaps.
\end{proof}

\begin{proposition}[Bounded-band genera]\label{prop:bounded-band-genera}
For every prime $p\leq A+1$ and $1\leq B\leq A$, each integer genus in the
interval
\[
  \left\lceil\frac B2\right\rceil(p-1)
  \leq g\leq B(p-1)
\]
is realized by a semigroup with $\AP(S)\leq A$.
\end{proposition}

\begin{proof}
Start with all $p-1$ coordinates equal to $L_B$.  Any integer increment from
$0$ to $(B-L_B)(p-1)$ can be distributed among the coordinates, each with
capacity $B-L_B$.  It follows that every intermediate coordinate sum occurs,
and the genus equals that sum.
\end{proof}

The bounded bands overlap as their upper coordinate bound varies.  Let
$p=P(A+1)$ and $n=p-1$.  For $1\leq B\leq A$, put
\[
  I_B=\left[\left\lceil\frac B2\right\rceil n,Bn\right].
\]
The intervals cover $[n,An]$: $I_1=\{n\}$, and
$\lceil(B+1)/2\rceil\leq B$ shows that consecutive intervals overlap.

\begin{theorem}[Realization up to the prime staircase]
\label{thm:realization-staircase}
For every $A\geq1$ and every integer $g$ satisfying
\[
  1\leq g\leq A(P(A+1)-1),
\]
there exists a numerical semigroup $S$ with
\[
  g(S)=g,
  \qquad \AP(S)\leq A.
\]
Equivalently, the set of realized genera satisfies
\[
  \boxed{\{g(S):\AP(S)\leq A\}
  \supseteq\{1,\ldots,A(P(A+1)-1)\}.}
\]
\end{theorem}

\begin{proof}
Proposition~\ref{prop:bounded-band-genera} and the interval coverage realize
every genus in $[n,An]$.  Since $n=p-1\leq A$, it remains only to realize
$1\leq g\leq A$.  The ordinary semigroup
$\{0,g+1,g+2,\ldots\}$ has gap set $\{1,\ldots,g\}$ and $\AP(S)=g$.
This proof is unconditional and does not use the prime staircase conjecture.
\end{proof}

This realization result allows the extremal function to be inverted at every
level where the staircase formula is known.  We define
\[
  a(g)=\min\{\AP(S):g(S)=g\},
  \qquad
  A_M(g)=\min\{A:M(A)\geq g\}.
\]
By definition, $A_M(g)\leq a(g)$.  A global equality does not follow formally:
if $M(A)$ were larger than the prime lower bound, the preceding theorem would
not realize the intervening genera.

\begin{corollary}[Inverse function at an exact staircase level]
\label{cor:inverse-exact-level}
Suppose that the exact staircase formula
\[
  M(A)=A(P(A+1)-1)
\]
holds.  It follows that every $1\leq g\leq M(A)$ is realized with
$\AP(S)\leq A$.  If
\[
  A_0=\min\{A:M(A)\geq g\}
\]
and the staircase formula is known at $A_0$, then
\[
  \boxed{a(g)=A_0.}
\]
\end{corollary}

\begin{proof}
Theorem~\ref{thm:realization-staircase} gives $a(g)\leq A_0$, and the
definition of $A_0$ gives $A_0\leq a(g)$.
\end{proof}

\begin{computercorollary}[Inverse staircase in the verified range]
\label{cor:verified-inverse}
For every $g\leq M(4\cdot10^{18})$, we have
\[
  \boxed{a(g)=\min\{A:g\leq M(A)\},}
\]
where in the relevant range
\[
  M(A)=
  \begin{cases}
    7,&A=3,\\
    A(P(A+1)-1),&A\neq 3.
  \end{cases}
\]
\end{computercorollary}

\begin{proof}
The minimizing $A$ is at most $4\cdot10^{18}$ by monotonicity.  Apply
Theorem~\ref{thm:verified-range} and
Corollary~\ref{cor:inverse-exact-level}; when $A=3$, genera through $6$ are
covered by Theorem~\ref{thm:realization-staircase}, and genus $7$ by
$\langle4,7,13\rangle$.
\end{proof}

Exact consecutive values of $M$ now translate directly into plateaux for the
inverse function.

\begin{corollary}\label{cor:inverse-plateaux}
Suppose the staircase formula is known at $A$ and $M(A-1)$ is known exactly.
Under these assumptions, the implication
\[
  M(A-1)<g\leq M(A)
  \quad\Longrightarrow\quad
  \boxed{a(g)=A.}
\]
holds.
In particular, this implication gives
\[
  \boxed{a(g)=10\qquad(55\leq g\leq100).}
\]
\end{corollary}

\begin{proof}
The lower bound follows from the definition of $M(A-1)$ and the upper bound
from Corollary~\ref{cor:inverse-exact-level}.  The example uses
$M(9)=54$ and $M(10)=100$; the latter follows from
Theorem~\ref{thm:square-characterization} because $11$ is prime.
\end{proof}

The plateau length is $M(A)-M(A-1)$ whenever the two consecutive values and
the realization at $A$ are known.  The set of levels at which the exact
formula holds simultaneously at $A-1$ and $A$ has density one, since the
exceptional set in Theorem~\ref{thm:density-one} and its unit translate both
have density zero.  This density statement concerns staircase levels and
does not by itself give a global closed formula for $a(g)$.

\section{Further directions}\label{sec:further}

The results above leave one main structural question unresolved: whether the
prime staircase formula holds for every \(A\neq 3\).  The prime-gap input
explains the finite verification and density-one results, but the harder
remaining step lies within the semigroup problem itself.  The arguments of
Sections~5 and~6 exclude most multiplicities using the loss from a single
divisor orbit; what remains is a narrow class of composite multiplicities for
which this first loss is insufficient.  Understanding how the different
orbit constraints interact with the Kunz inequalities is the central task.

Two secondary directions concern settings in which the invariant may be
useful beyond the extremal problem itself.  Weierstrass semigroups and
one-point algebraic-geometry codes provide one such setting, because the
arrangement of the gaps already plays an important role.  Another consists of
special classes of numerical semigroups, such as symmetric semigroups, for
which arithmetic progressions in the gap set can be related directly to
arithmetic progressions among the nongaps.

To isolate the remaining structural obstruction, let \(p=P(A+1)\), and
suppose that a composite multiplicity \(m>p\) survives the single-divisor
criterion of Theorem~\ref{thm:prime-gap-criterion}.
Lemma~\ref{lem:least-factor-loss} then forces
\[
    m-p>\sqrt m-2.
\]
Thus, a surviving multiplicity can occur only when \(A+1\) lies unusually far
from the preceding prime.  Moreover, the least-factor estimate is weakest when
the least prime factor of \(m\) is close to \(\sqrt m\).  This singles out prime
squares and balanced semiprimes as the most natural cases to understand.

The first example already occurs at \(A=9\).  The multiplicities
\(9=3^2\) and \(10=2\cdot 5\) lie between the consecutive primes \(7\) and
\(11\), and the single-divisor estimate does not recover the sharp bound.
For \(m=9\), the additional information comes from anchoring.  The Run Lemma
controls not only the minimum Kunz coordinate on an orbit but also the
position at which that minimum first occurs.  Once this positional information
is combined with the Kunz inequalities, bounds propagate to further
coordinates and recover the correct genus.

The same mechanism should persist more generally.  For instance, when
\(m=q^2\), every nonzero \(q\)-orbit satisfies
\[
    \min_{i\in\mathcal O} k_i\leq
    \left\lfloor\frac{A}{q}\right\rfloor .
\]
These small coordinates cannot be placed independently because the Kunz
inequalities relate them through addition in \(\Z/q^2\Z\).
The problem is therefore not only to measure the loss on each orbit but also to
understand how minima on different orbits force one another.

This suggests the following more precise question.

\begin{problem}[Propagated loss]\label{prob:propagated-loss}
Find a quantitative genus-loss estimate for prime-square and balanced
semiprime multiplicities that combines Run-Lemma anchoring with Kunz
subadditivity, and that is strong enough to rule out the remaining composite
competitors in the prime staircase conjecture.
\end{problem}

Even a sufficiently uniform additional loss would already have consequences.
For example, a bound of the form
\[
    g\leq A(m-1)-cAm
\]
with some fixed \(c>0\) in the residual cases, together with
\(P(x)=(1-o(1))x\), would eliminate such multiplicities for all sufficiently
large \(A\).  A sharper analysis of the propagation phenomenon may therefore
lead to a proof of the prime staircase conjecture without any hypothesis on
prime gaps.

Beyond the extremal problem, Weierstrass semigroups provide a natural setting
in which the arrangement of gaps matters.  If \(P\) is a rational point on a
smooth projective curve, then its Weierstrass semigroup \(H(P)\) is numerical
and its gap set has cardinality equal to the genus
\cite{stichtenoth-function-fields}.  In the theory of one-point
algebraic-geometry codes, the Feng--Rao order bound depends on additive
relations among nongaps, while consecutive Weierstrass gaps are already known
to influence distance estimates
\cite{stichtenoth-function-fields,garcia-kim-lax}.  It is therefore natural to
ask whether the finer spectrum \(\Sigma(H(P))\), rather than only
\(\operatorname{AP}(H(P))\), carries useful information about the Feng--Rao
sequence or the order bound.  Structured Weierstrass semigroups, particularly
those arising from recursive or otherwise highly organized families, provide
natural test cases for this question.

Symmetric numerical semigroups provide another structured class.  In this
case, arithmetic progressions of gaps admit a simple reflection across the
Frobenius number.  Suppose that
\[
x,x+d,\ldots,x+(k-1)d\subseteq G(S).
\]
Symmetry then gives the reflected progression
\[
F-x-(k-1)d,\ldots,F-x\subseteq S\cap[0,F].
\]
Consequently, the part of the spectrum \(\Sigma(S)\) arising from gap
differences \(d\in G(S)\setminus\{1\}\) may also be viewed through arithmetic
progressions among the nongaps below \(F\).  It would be interesting to
determine how much of the gap spectrum can be recovered from this reflected
picture.

These considerations lead to four concrete questions.

\begin{enumerate}
    \item Establish a propagated genus-loss theorem for the residual
    prime-square and balanced-semiprime multiplicities by combining
    Run-Lemma anchoring with Kunz subadditivity.  Such a theorem is the main
    structural route to proving the prime staircase conjecture for every
    \(A\neq3\) without prime-gap hypotheses or finite computation.

    \item Find a conceptual replacement for the remaining finite certificates,
    especially the multiplicities \(9\) and \(10\) at \(A=9\).

    \item Determine whether \(A=3\) is the only exception to the prime
    staircase formula and, separately, classify the semigroups attaining
    \(M(A)\).  The value predicted by the conjecture does not imply uniqueness
    of a maximizing semigroup.

    \item Understand the gap-difference part of \(\Sigma(S)\), in particular
    for structured classes such as symmetric and Weierstrass semigroups.
\end{enumerate}

The first question is the immediate continuation of the present work.  The
calculation at \(A=9\) suggests that the missing ingredient is a systematic
way of propagating the positional information supplied by the Run Lemma
through the Kunz inequalities.

\appendix

\section{Computational certificates and reproducibility}
\label{app:computation}

All filenames in this appendix refer to files in the accompanying directory
\compdir{}.

The finite search used in Proposition~\ref{prop:finite-certificate} was
performed by \path{plateau_certificate.c}; its corresponding output is
designated \path{plateau_certificate.log}.  It enumerates the bounded Kunz boxes for
$m=8,9,10$ and uses only the Kunz inequalities and the explicitly stated
Run-Lemma constraints.  It does not construct full gap sets or compute
$\AP_d(S)$ directly, and it does not assert that these necessary constraints
characterize the condition $\AP(S)\leq9$.  A second exhaustive enumeration,
implemented in \path{plateau_certificate_independent.c}, uses a different
branching order and a disjunctive encoding of the orbit conditions.  It
produced the same maxima and the same unique maximizing vectors; its output is
recorded in \path{plateau_certificate_independent.log}.  Both implementations
and their outputs accompany this draft.

As an independent check, the program \path{direct_ap_search.c}
constructs gap sets and computes their arithmetic progressions directly.
Its \texttt{a8} mode proves Theorem~\ref{thm:M8-computation}; the output is
designated \path{direct_ap_A8.log}.  Its \texttt{table} mode produces the
comparison displayed after Conjecture~\ref{conj:prime-staircase}; the output
is designated \path{direct_ap_table_12.log}.  Unlike the finite Kunz
certificate, this program constructs each candidate gap set and computes the
arithmetic progressions themselves.  Its \texttt{extremizer9} mode directly
checks the multiplicity-$9$ sharpness example, with output in
\path{direct_ap_extremizer9.log}, while \texttt{tests} runs the elementary
regression checks recorded in \path{direct_ap_tests.log}.

The exact inequality calculation in
Proposition~\ref{prop:criterion-computation} was performed by
\path{staircase_criterion.c}; the file
\path{staircase_criterion_1e8.log} records the run through $10^8$.  This
calculation factors the relevant multiplicities and tests the exact composite
criterion; it does not enumerate numerical semigroups.  The
companion \path{staircase_chunk.c} performs the same calculation on a
specified interval, with a sample output in \path{staircase_chunk_sample.log}.
A second complete execution on $1\leq A\leq10^8$ is recorded in
\path{staircase_chunk_1e8.log}; it returns the same pair count and the same
four failures as \path{staircase_criterion_1e8.log}.

For reproducibility, the direct search uses dynamically sized Kunz arrays and
a dynamically allocated gap array indexed from $0$ through the computed
Frobenius number.
An earlier fixed-width 128-bit mask overflowed when a possible gap position
exceeded $127$, producing a false discrepancy at $A=12$; widening the
representation removed it.  Future extensions should allocate on the order
of $A(A+1)$ positions and avoid fixed Kunz arrays tied to a small
multiplicity.

% TODO before submission: reproduce the 10^8 criterion computation
% independently, preferably with
% computational_reproducibility_files/staircase_chunk.c on a different compiler
% and machine, and retain checksums and complete log files.

\section*{Acknowledgments}

An early version of this project began to take shape during AGRA V---A School on
Arithmetic, Groups and Analysis, held in Popayán, Colombia, from July 27 to
August 7, 2026.  I am grateful to the organizers for providing a stimulating
environment in which some of the first ideas and computations behind this work
were developed.

\section*{Funding}

The author was supported by the Fundação Carlos Chagas Filho de Amparo à Pesquisa
do Estado do Rio de Janeiro (FAPERJ) through the Nota 10 Excellence Scholarship.

\section*{Conflict of interest}

The author declares no conflict of interest.

\section*{Code and data availability}

The source code, complete computational logs, concise certificate, checksums,
and reproduction instructions are publicly available at
\url{https://github.com/alequisGS/alequisGS.github.io/tree/main/drafts/arithmetic-progressions-gap-sets/computations}.
The repository contains the primary and independent plateau certificates, the
direct AP checker, the two staircase-criterion programs, and their complete
logs.  Compilation and reproduction commands appear in its
\path{README.md}; the generated summary and checksum manifest are
\path{CERTIFICATE.txt} and \path{SHA256SUMS}.  A versioned archival release
with a permanent DOI will be deposited on Zenodo before publication.

% TODO BEFORE PUBLICATION: archive a versioned release on Zenodo and add its
% permanent DOI to the preceding paragraph.

\printbibliography

\end{document}
